%Author:Haoyu Tang 
%Github:Jerry-Haoyu 
%Date:2026
%This is a free-for-all template, cheers LATEX nerds! 

% For more stylish math
\documentclass[14pt]{extarticle}
\usepackage{unicode-math}
\setmathfont{STIX Two Math}

%font
\RequirePackage[scaled]{berasans}
\RequirePackage[semibold]{sourcecodepro}
\RequirePackage[scaled=.96,osf,sups]{XCharter}


\usepackage{graphicx} % Required for inserting images
\usepackage[dvipsnames]{xcolor}
\usepackage{amsmath}
\usepackage{cancel}
\usepackage{amssymb}
\usepackage{amsthm}
\usepackage{tikz}
\usepackage[most]{tcolorbox}
\usepackage{geometry}
 \geometry{
 a4paper,
 total={170mm,300mm},
 left=25mm,
 right=25mm,
 top=26mm,
 bottom=26mm
 }

%set cover, with hyperlinked table of content
\newcommand{\cover}[4]{
    \begin{titlepage}
        \centering
        \vspace*{6cm} 
        {\LARGE\bfseries  {#1} \par} % Title
        \vspace{2.5cm}
        {\LARGE {#2} \par}
        \vspace{1cm}
        {\Large #3\par} % Author
        \vspace{3cm}
        {\large \today\par} % Date
        \vfill 
        {\small #4\par}
        \newpage 
        \tableofcontents
    \end{titlepage}
 }

%set header using fancyhdr
\newcommand{\setheader}[3]{
    \usepackage{fancyhdr}
    \pagestyle{fancy}
    \fancyhf{} % clear all header and footer fields
    \fancyfoot[L]{\thedate}
    \fancyhead[L]{#1}
    \fancyhead[C]{#2}
    \fancyhead[R]{#3}
}


%===========================
%      Some Scafolding
%===========================

%colored version of various sections title
\newcommand{\bsection}[1]{\section{\textcolor{RoyalBlue}{#1}}}
\newcommand{\bsubsection}[1]{\subsection{\textcolor{RoyalBlue}{#1}}}
\newcommand{\bsubsubsection}[1]{\subsubsection{\textcolor{RoyalBlue}{#1}}}

\newcommand{\bsectionnn}[1]{\section*{\textcolor{RoyalBlue}{#1}}}
\newcommand{\bsubsectionnn}[1]{\subsection*{\textcolor{RoyalBlue}{#1}}}
\newcommand{\bsubsubsectionnn}[1]{\subsubsection*{\textcolor{RoyalBlue}{#1}}}



\newenvironment{tbox}{
    \begin{tcolorbox}[colframe=white,sharp corners]
    
}{
    \end{tcolorbox}
}


\newenvironment{pbd}{
    \begin{tcolorbox}[colframe=white,sharp corners]
    \noindent \textbf{Problem Description:}
}{
    \end{tcolorbox}
}

\newenvironment{rmk}{
    \textbf{Rmk.}
}{
    \newline
}

\newenvironment{sol}{
    \newline
    \noindent \textbf{Solution.}
}{
    \newline
}




%===========================
%      Some Basic Math
%===========================
% \theoremstyle{definition} %for up right text
\newtheorem{prop}{Prop.}
\newtheorem{ex}{Ex.}
\newtheorem{df}{Def.}
\newtheorem{thm}{Thm.}
\newtheorem{coro}{Corollary.}
\newtheorem{lemma}{Lemma.}

\newcommand{\rn}[1]{\mathbb{R}^{#1}}
\newcommand{\cn}[1]{\mathbb{C}^{#1}}
\newcommand{\z}{\mathbb{Z}}

\newcommand{\argmin}[1]{\mathrm{argmin}_{#1}}
\newcommand{\argmax}[1]{\mathrm{argmax}_{#1}}
\newcommand{\ind}[1]{\mathbb{I}\{#1\}}

%some probability
\newcommand{\E}[1]{\mathbb{E}\left[#1\right]}
\renewcommand{\Pr}[1]{\mathrm{Pr}\left(#1\right)}

%some vector calculus 

\newcommand{\dvg}[1]{(\nabla \cdot #1)}
\newcommand{\crl}[1]{\nabla  \times #1}

\renewcommand{\vec}[1]{\mathbf{#1}}


%===========================
%      Psuedo Code
%===========================

%previous version exploting some turing-compelteness LOL
%command for psuedocode
% \newcommand{\algo}[1]{
%     \textup{
%         \begin{tabular}{ll}
%             #1 
%         \end{tabular}
%     }
% }
% \newcommand{\class}[2]{
%     \texttt{class \; #1} \; \{ & \\
%     #2 
%     \} & \\
% }
% \newcommand{\tab}{
% \qquad
% }
% \newcommand{\code}[1] {
% \texttt{#1}  & \\
% }

% \newcommand{\cmt}[1] {
% \textcolor{gray}{#1} & \\
% }

% \newcommand{\function}[4] {
%     \texttt{#1 \; #2(#3)} \; \{ & \\
%         #4
%     \} & \\
% }

% \newcommand{\for}[2] {
%      \texttt{for #1} \; \{ & \\
%         #2
%      \} & \\
% }
% \newcommand{\indfunction}[4] {
%      \tab \texttt{#1 #2(#3)} \; \{ & \\
%         #4
%      \tab \} & \\
% }

% \newcommand{\indfor}[2] {
%      \tab \texttt{for #1} \; \{ & \\
%         #2
%      \tab \} & \\
% }


%TYPE 1: for real code and lined psuedo 
\usepackage{listings}

\definecolor{codegreen}{rgb}{0,0.6,0}
\definecolor{codegray}{rgb}{0.5,0.5,0.5}
\definecolor{codepurple}{rgb}{0.58,0,0.82}
\definecolor{backcolour}{rgb}{0.95,0.95,0.92}

\lstdefinestyle{mystyle}{
    backgroundcolor=\color{white!95!black},   
    commentstyle=\color{codegreen},
    keywordstyle=\color{magenta},
    numberstyle=\tiny\color{codegray},
    stringstyle=\color{codepurple},
    basicstyle=\ttfamily\footnotesize,
    breakatwhitespace=false,         
    % breaklines=true,                 
    captionpos=b,                    
    keepspaces=true,                 
    numbers=left,                    
    numbersep=5pt,                  
    showspaces=false,                
    showstringspaces=false,
    showtabs=false,                  
    tabsize=2
}
\lstset{style=mystyle}


%type 2: for psuedo code using algorithmic
\usepackage{algorithmic}
%defining tcolorbox algobox
\newtcolorbox{algobox}[2][]{%
  enhanced,
  breakable,
  colback=white,
  colframe=black,
  fonttitle=\bfseries,
  title={Algorithm: #2},
  sharp corners,
  width=0.85\textwidth,  % Adjust width here
  boxrule=0.8pt,
  #1
}
%===========================
%      For hyperlinks
%===========================
\usepackage{hyperref}
\hypersetup{
    colorlinks=true,
    linkcolor=blue,
    filecolor=magenta,      
    urlcolor=cyan,
    pdftitle={Overleaf Example},
    pdfpagemode=FullScreen,
}

%===========================
%          *END*
%===========================



\cover{Realistic Simulation of Galewsky Initial Value Problem}{Haoyu Tang}{}{}

\begin{document}

\section{Psuedospectral Method}
There are four equivalent ways of describing the SWE system on the sphere:
\begin{figure}[h]
    \centering
    

\tikzset{every picture/.style={line width=0.75pt}} %set default line width to 0.75pt        

\begin{tikzpicture}[x=0.75pt,y=0.75pt,yscale=-1,xscale=1]
%uncomment if require: \path (0,300); %set diagram left start at 0, and has height of 300

%Left Right Arrow [id:dp2543051183422924] 
\draw  [color={rgb, 255:red, 0; green, 0; blue, 0 }  ,draw opacity=0.05 ][fill={rgb, 255:red, 194; green, 158; blue, 249 }  ,fill opacity=1 ] (162.91,50.1) -- (182.4,78.72) -- (172.72,78.82) -- (172.97,136.44) -- (182.65,136.34) -- (163.41,165.36) -- (143.92,136.74) -- (153.61,136.64) -- (153.35,79.02) -- (143.67,79.12) -- cycle ;
%Left Right Arrow [id:dp7369355028476874] 
\draw  [color={rgb, 255:red, 0; green, 0; blue, 0 }  ,draw opacity=0.1 ][fill={rgb, 255:red, 74; green, 144; blue, 226 }  ,fill opacity=0.3 ] (202.86,49.81) -- (211.59,78.54) -- (207.28,78.58) -- (207.54,136.21) -- (211.84,136.17) -- (203.36,165.07) -- (194.63,136.34) -- (198.93,136.3) -- (198.68,78.67) -- (194.38,78.72) -- cycle ;
%Left Right Arrow [id:dp47994436496181625] 
\draw  [color={rgb, 255:red, 0; green, 0; blue, 0 }  ,draw opacity=0.05 ][fill={rgb, 255:red, 194; green, 158; blue, 249 }  ,fill opacity=1 ] (406.55,45.95) -- (426.04,74.56) -- (416.36,74.66) -- (416.61,132.29) -- (426.3,132.19) -- (407.06,161.2) -- (387.57,132.59) -- (397.25,132.49) -- (397,74.86) -- (387.32,74.96) -- cycle ;
%Left Right Arrow [id:dp19367934898724093] 
\draw  [color={rgb, 255:red, 0; green, 0; blue, 0 }  ,draw opacity=0.1 ][fill={rgb, 255:red, 74; green, 144; blue, 226 }  ,fill opacity=0.3 ] (449.61,45.66) -- (458.34,74.38) -- (454.04,74.43) -- (454.29,132.05) -- (458.59,132.01) -- (450.12,160.91) -- (441.39,132.19) -- (445.69,132.14) -- (445.44,74.52) -- (441.14,74.56) -- cycle ;
%Curve Lines [id:da22359335462503038] 
\draw    (170.6,185.6) .. controls (205.88,162.56) and (371.38,157.93) .. (409.48,183.81) ;
\draw [shift={(410.6,184.6)}, rotate = 216.72] [color={rgb, 255:red, 0; green, 0; blue, 0 }  ][line width=0.75]    (10.93,-3.29) .. controls (6.95,-1.4) and (3.31,-0.3) .. (0,0) .. controls (3.31,0.3) and (6.95,1.4) .. (10.93,3.29)   ;
%Curve Lines [id:da8324860879521677] 
\draw    (410.6,204.4) .. controls (375.95,236.87) and (210.95,236.61) .. (173.69,211.37) ;
\draw [shift={(172.6,210.6)}, rotate = 36.61] [color={rgb, 255:red, 0; green, 0; blue, 0 }  ][line width=0.75]    (10.93,-3.29) .. controls (6.95,-1.4) and (3.31,-0.3) .. (0,0) .. controls (3.31,0.3) and (6.95,1.4) .. (10.93,3.29)   ;
%Shape: Rectangle [id:dp8400991491232823] 
\draw  [color={rgb, 255:red, 0; green, 0; blue, 0 }  ,draw opacity=0 ][fill={rgb, 255:red, 155; green, 155; blue, 155 }  ,fill opacity=0.17 ] (377.6,168.2) -- (465.6,168.2) -- (465.6,216.8) -- (377.6,216.8) -- cycle ;

% Text Node
\draw (153.28,14.44) node [anchor=north west][inner sep=0.75pt]    {$( u,v,\Phi )$};
% Text Node
\draw (150.17,182.86) node [anchor=north west][inner sep=0.75pt]    {$(\hat{u} ,\hat{v} ,\hat{\Phi })$};
% Text Node
\draw (389,16.52) node [anchor=north west][inner sep=0.75pt]    {$( \zeta ,\ \delta ,\ \Phi )$};
% Text Node
\draw (389,180.78) node [anchor=north west][inner sep=0.75pt]    {$(\hat{\zeta } ,\ \hat{\delta } ,\ \hat{\Phi })$};
% Text Node
\draw (263.69,121.04) node [anchor=north west][inner sep=0.75pt]    {$-\frac{l( l+1)}{a}$};
% Text Node
\draw (261.69,232.04) node [anchor=north west][inner sep=0.75pt]    {$-\frac{a}{l( l+1)}$};
% Text Node
\draw (153.21,127.51) node [anchor=north west][inner sep=0.75pt]  [rotate=-269.93]  {$vSHT$};
% Text Node
\draw (396.21,122.51) node [anchor=north west][inner sep=0.75pt]  [rotate=-269.93]  {$vSHT$};
% Text Node
\draw (194.21,125.51) node [anchor=north west][inner sep=0.75pt]  [rotate=-269.93]  {$sSHT$};
% Text Node
\draw (441.21,120.51) node [anchor=north west][inner sep=0.75pt]  [rotate=-269.93]  {$sSHT$};


\end{tikzpicture}

\end{figure}

The transformation between $(u,v, \Phi)\leftrightarrow (\hat u, \hat v, \hat \Phi)$ and similarly for $(\zeta, \delta)$
is the real spherical harmonic transform. In particular, for $(u,v)\leftrightarrow (\hat u, \hat v)$, we use the \textbf{real vector spherical harmonic transform}
defined as:

\begin{tbox}
    \begin{df}
        Let $\mathbf{v}$ be a vector field $S^2\to \mathbb{R}^2$. Consider the multipole expansion:
        \begin{equation}
            \mathbf{v} = \sum_{l=1}^{\infty} \sum_{m=-l}^l \psi_{lm}\boldsymbol{\Psi}_{lm} +  \phi_{lm}\boldsymbol{\Phi}_{lm}
        \end{equation}
        where 
        $$
        \begin{cases}
            \Psi_{lm} &= \nabla Y_{lm}(\theta,\phi) \\ 
            \Phi_{lm} &= \vec k\times \nabla Y_{lm}(\theta,\phi)
        \end{cases}
        $$
        where $\Psi_{lm}$ is the curl-free \textbf{Poloidal Harmonics}.
        $\Phi_{lm}$ is the divergence-free \textbf{Toroidal Harmonics}.
        $\vec k$ is the unit vector pointing outward from the surface.
        Define the vector spherical harmonics transform $\mathcal F_{vSHT}$ as:
        \begin{align}
            \mathcal F_{vSHT} : L^2(TS^2) \to C(A;\mathbb{C})
        \end{align}
        where $TS^2$ is the tangent bundle on $S^2$ and $A=\{(l,m):l\in \mathbb{N}^+,-l\leq m\leq l\}$
    \end{df} 
\end{tbox}

Some useful properties of vector spherical harmonics:
\begin{tbox}
    Azimuthal symmetry up to sign and complex conjugation:
    \begin{prop}
        \begin{align}
            \Psi_{l,-m} &= (-1)^m \Psi_{l,m}^* \\
            \Phi_{l,-m} &= (-1)^m \Phi_{l,m}^*
        \end{align}
    \end{prop}
    Orthognality:
    \begin{prop}
        \begin{align}
            \langle \Psi_{lm}, \Psi_{l'm'}\rangle  &= l(l+1)\delta_{ll'}\delta_{mm'} \\
            \langle \Phi_{lm}, \Phi_{l'm'}\rangle  &= l(l+1)\delta_{ll'}\delta_{mm'} \\
            \langle \Psi_{lm}, \Phi_{l'm'} \rangle &= 0 
        \end{align}
    \end{prop}
\end{tbox}
The first property allows us to cleanly define the truncated version of vSHT:
\begin{tbox}
    \begin{df}
        Define the vector spherical harmonics transform as: 
        \begin{align}
            \mathrm{vSHT^{n_g,l_{max},m_{max}}}: n_g\times 2n_g &\to l_{max}\times m_{max} \\
            \{(u_i, v_j)\}& \mapsto \{\phi_{lm},\psi_{lm}\}
        \end{align}
        where $n_g$ defines the grid size corresponding to latitude and $l_{max}, m_{max}$ specifies the truncation 
        of the multipole expansion.
    \end{df}
\end{tbox}
The second property gives an easy transformation between $(\hat u, \hat v)\leftrightarrow (\hat \zeta, \hat \delta)$, recall:
\begin{tbox}
    \begin{prop}
        $$(\hat \zeta, \hat \delta) = -\frac{l(l+1)}{a}\mathcal{F}_{vSHT}(u,v)$$
        With the inverse transformation given by:
        $$(\hat u, \hat v) = -\frac{a}{l(l+1)}\mathcal{F}_{vSHT}(\hat \zeta, \hat \delta)$$
    \end{prop}
\end{tbox}
by applying the definition $\zeta = k \cdot \nabla \times \mathbf{v}$, $\delta = \nabla \cdot \mathbf{v}$ to the multipole expansion of 
$\mathbf{v}$. 
\end{document}